\SourcePage{42}{47}
\section{The two-photon system}

Arguments analogous to those developed in \secref{6} make it possible to count the number of allowed states in the more involved case of a system of two photons as well (\emph{L.~Landau}, 1948).

Let us work in the center-of-inertia frame of the two photons, where $\mathbf{k}_1 = -\mathbf{k}_2 \equiv \mathbf{k}$\,\footnote{A center-of-inertia frame can be found in every case except when the two photons propagate in the same direction. For such a pair the total momentum $\mathbf{k}_1 + \mathbf{k}_2$ and total energy $\omega_1 + \omega_2$ are connected by exactly the same relation as for one photon, so no frame with $\mathbf{k}_1 + \mathbf{k}_2 = 0$ exists.}. In the momentum representation the wave function of the pair can be cast as a three-dimensional rank-2 tensor $A_{ik}(\mathbf{n})$ built bilinearly from the components of the individual photon vector wave functions; each of its two indices refers to one photon ($\mathbf{n}$ denotes the unit vector in the direction of~$\mathbf{k}$). The transversality requirement for each photon is encoded in the orthogonality of~$A_{ik}$ to~$\mathbf{n}$:
\begin{equation}
A_{il}n_l = 0, \qquad A_{lk}n_l = 0.
\tag{9.1}
\end{equation}

Interchanging the two photons amounts to transposing the indices of the tensor~$A_{ik}$ together with a simultaneous reversal of the sign of~$\mathbf{n}$. Because photons obey Bose statistics,
\begin{equation}
A_{ik}(-\mathbf{n}) = A_{ki}(\mathbf{n}).
\tag{9.2}
\end{equation}

The tensor~$A_{ik}$ is, in general, not symmetric in its indices. We decompose it into its symmetric part~($s_{ik}$) and antisymmetric part~($a_{ik}$): $A_{ik} = s_{ik} + a_{ik}$. Each of these parts must individually satisfy the relation~(9.2) (as well as the orthogonality conditions~(9.1)). This gives
\begin{equation}
s_{ik}(-\mathbf{n}) = s_{ik}(\mathbf{n}),
\tag{9.3}
\end{equation}
\begin{equation}
a_{ik}(-\mathbf{n}) = -a_{ik}(\mathbf{n}).
\tag{9.4}
\end{equation}

Spatial inversion does not by itself reverse the sign of the components of a second-rank tensor, but it does reverse the sign of~$\mathbf{n}$. From~(9.3) it is therefore apparent that the wave function~$s_{ik}$ is symmetric under inversion, i.e.\ it corresponds to even-parity states of the photon system; the wave function~$a_{ik}$, on the other hand, corresponds to odd-parity states.
\SourcePage{43}{48}

An antisymmetric second-rank tensor is equivalent (dual) to an axial vector~$\mathbf{a}$ whose components are expressed via the tensor components as $a_i = \tfrac{1}{2}e_{ikl}a_{kl}$, where~$e_{ikl}$ is the antisymmetric unit tensor (see~II, \S~6). The orthogonality of the tensor~$a_{kl}$ to the vector~$\mathbf{n}$ implies that the vectors~$\mathbf{a}$ and~$\mathbf{n}$ are parallel\,\footnote{We have: $a_{ik} = e_{ikl}a_l$, and the orthogonality condition yields
\[
a_{ik}n_k = e_{ikl}a_ln_k = [\mathbf{na}]_i = 0.
\]}. One may therefore write $\mathbf{a} = \mathbf{n}\varphi(\mathbf{n})$, where~$\varphi$ is a scalar; by~(9.4) we must have $\mathbf{a}(-\mathbf{n}) = -\mathbf{a}(\mathbf{n})$, and hence
\[
\varphi(-\mathbf{n}) = \varphi(\mathbf{n}).
\]

This relation means that the scalar~$\varphi$ can be built linearly only from spherical harmonics of even order~$L$ (including zero).

We see that the antisymmetric tensor~$a_{ik}$, by its transformation properties (under rotations), is equivalent to a single scalar (cf.\ the footnote on p.~34). Assigning to the latter a ``spin'' of~0, we find that the angular momentum of the state is $J = L$. Thus the tensor~$a_{ik}$ corresponds to odd-parity states of the photon system with even total angular momentum~$J$.

Let us turn to the symmetric tensor~$s_{ik}$. Since it is even under reversal of the sign of~$\mathbf{n}$, it describes even-parity states of the photon system. It follows likewise that all components of~$s_{ik}$ are expressed through spherical harmonics of even order~$L$ (including $L = 0$). An arbitrary symmetric second-rank tensor can be decomposed, as is well known, into a scalar~($s_{ii}$) and a traceless symmetric tensor~($s'_{ik}$) with $s'_{ii} = 0$.

The scalar~$s_{ii}$ is assigned ``spin''~0, so the angular momentum of the corresponding states is $J = L$, which is even. The tensor~$s'_{ik}$ corresponds to ``spin''~2 (see~III, \S~57). Composing this ``spin'' with the even ``orbital angular momentum''~$L$ according to the addition rule for angular momenta, we find that for a given even $J \neq 0$ there are three states (with $L = J \pm 2, J$), while for odd $J \neq 1$ there are two states (with $L = J \pm 1$). The exceptions are $J = 0$ with a single state ($L = 2$) and $J = 1$ with a single state ($L = 2$).

In this counting, however, we have not yet accounted for the orthogonality of the tensor~$s_{ik}$ to the vector~$\mathbf{n}$. One must therefore subtract from the number of states obtained those corresponding to a symmetric second-rank tensor that is ``parallel'' to the vector~$\mathbf{n}$. Such a tensor (denote it~$s''_{ik}$) can be written as
\[
s''_{ik} = n_ib_k + n_kb_i,
\]
where~$\mathbf{b}$ is some vector. By~(9.3) this vector must satisfy $\mathbf{b}(-\mathbf{n}) = -\mathbf{b}(\mathbf{n})$. The tensor~$s''_{ik}$ responsible for the ``extra'' states is thus equivalent to an odd-parity vector. This vector must consequently be expressed through spherical harmonics of odd orders~$L$ only. Noting
\SourcePage{44}{49}
further that a vector corresponds to ``spin''~1, we find that for each even $J \neq 0$ there are two states (with $L = J \pm 1$), and for each odd~$J$ a single state (with $L = J$); the special case $J = 0$ has one state (with $L = 1$).

Collecting all the results obtained, we arrive at the following table giving the number of possible even- and odd-parity states of a system of two photons (with vanishing total momentum) for the various values of the total angular momentum~$J$:
\begin{equation}
\begin{array}{c|c|c}
J & \text{Even} & \text{Odd} \\ \hline
0 & 1 & 1 \\
1 & - & - \\
2k & 2 & 1 \\
2k+1 & 1 & -
\end{array}
\tag{9.5}
\end{equation}
($k$ being a positive integer other than zero). We observe that for odd~$J$ no odd-parity states exist, while the value $J = 1$ is altogether excluded\,\footnote{For an alternative derivation of these results, see problem~1 to \secref{69}.}.

The wave function of the two-photon system~$A_{ik}$ determines the correlation between their polarizations. The probability that both photons simultaneously possess definite polarizations~$\mathbf{e}_1$ and~$\mathbf{e}_2$ is proportional to
\[
A_{ik}e_{1i}^*e_{2k}^*.
\]
In other words, if the polarization~$\mathbf{e}_1$ of one photon is specified, then the polarization of the second photon,~$\mathbf{e}_2$, is proportional to
\begin{equation}
e_{2k} \propto A_{ik}e_{1i}^*.
\tag{9.6}
\end{equation}
In odd-parity states of the system, $A_{ik}$ coincides with the antisymmetric tensor~$a_{ik}$. In that case
\[
\mathbf{e}_2\mathbf{e}_1^* \propto a_{ik}e_{1i}^*e_{1k}^* = 0,
\]
i.e.\ the polarizations of the two photons are mutually orthogonal. For linear polarization this signifies that their directions are perpendicular, while for circular polarization it means opposite senses of rotation.

The even-parity state with $J = 0$ is described by a symmetric tensor reducing to a scalar,
\[
s_{ik} = \text{const} \cdot (\delta_{ik} - n_in_k).
\]
From~(9.6) we then obtain $\mathbf{e}_1 = \mathbf{e}_2^*$. For linear polarization this means that their directions are parallel, whereas for circular polarization it means opposite senses of rotation. The latter fact is self-evident: when $J = 0$ the sum of the angular-momentum projections of the two photons onto any common direction~$\mathbf{k}$ must in all cases vanish (the projections onto the opposite directions~$\mathbf{k}_1$ and~$\mathbf{k}_2$, i.e.\ the helicities, are accordingly equal).

